\begin{helper}
Assume \(\noderef{FiberAndDegreeEvent}\).  For any finite set \(I\) and any
two distinct base vertices \(x\ne y\),
\[
|X_x(I)\cap X_y(I)|\le 100(\log n)^3.
\]
\end{helper}
\begin{proof}
Unfold \(X_z(I)\).  It is the subset of \(I\) consisting of vertices in the
lifted neighborhood of \(z\).  Therefore
\[
X_x(I)\cap X_y(I)\subseteq N_x\cap N_y,
\]
where \(N_z\) denotes the lifted neighborhood of the base vertex \(z\).

Write the two base vertices in the disjoint union of red and blue base
vertices.  If both are red, then \(x\ne y\) is exactly the distinctness of the
two red indices, and the red-red lifted-neighborhood intersection clause of
\(\noderef{FiberAndDegreeEvent}\) gives the bound for \(N_x\cap N_y\).  If
both are blue, the blue-blue clause gives the same bound.  If one is red and
one is blue, the mixed red-blue clause gives the same bound; reversing the
order of the two colors only uses commutativity of intersection.  In every
case, finite-set cardinality monotonicity along the displayed inclusion gives
the asserted bound for \(X_x(I)\cap X_y(I)\).
\end{proof}
